\section*{الفصل \ref{ch:g3:problems} --- حلّ المسائل}

\begin{solution}{exo:g3:problems:1}
انضمام القطيعين: جمع. وتقسيم الكرز: قسمة. وعلب
الأقلام: ضرب. وإعطاء البطاقات: طرح.
\end{solution}

\begin{solution}{exo:g3:problems:2}
المطلوب: عدد السيارات الآن. والمعلوم: $246$، ثم $158$ أخرى.
والعملية: جمع. والحساب: $246 + 158 = 404$. والجملة: في الموقف
الآن $404$ سيارة.
\end{solution}

\begin{solution}{exo:g3:problems:3}
قسمة تجميع: $90 \div 9 = 10$. فيعطي الحبل $10$ قطع.
\end{solution}

\begin{solution}{exo:g3:problems:4}
ضرب: $26 \times 8 = 208$. فتكلّف التذاكر $208$.
\end{solution}

\begin{solution}{exo:g3:problems:5}
شريط مقارنة: شريط هدى $64$، وشريط هشام أقصر بمقدار $29$.
والطرح: $64 - 29 = 35$. فعند هشام $35$ ملصقًا.
\end{solution}

\begin{solution}{exo:g3:problems:6}
المستلَم: $6 \times 45 = 270$ تفاحة. \omterm{def:g3:division:remainder}{والباقي}:
$270 - 178 = 92$ تفاحة.
\end{solution}

\begin{solution}{exo:g3:problems:7}
كتب العلوم: $6 + 8 = 14$. والروايات: $12 - 9 = 3$ أكثر في الأسبوع 2.
\end{solution}

\begin{solution}{exo:g3:problems:8}
الاثنين $12$، والثلاثاء $20$، والأربعاء $12$: \omterm{def:g1:addition:def}{والمجموع}
$12 + 20 + 12 = 44$ كعكة. وأفضل يوم هو الثلاثاء.
\end{solution}

\begin{solution}{exo:g3:problems:9}
العدد الذي لا فائدة منه هو الثمن ($12$). والقسمة:
$210 \div 35 = 6$ (لأن $35 \times 6 = 210$). فيلزم لؤيًا $6$ أيام.
\end{solution}

\begin{solution}{exo:g3:problems:10}
الكبار: $2 \times 9 = 18$. والأطفال: $3 \times 4 = 12$. والثمن الكلي:
$18 + 12 = 30$. \omterm{def:g3:division:remainder}{والباقي}: $50 - 30 = 20$. فيستردّون $20$.
\end{solution}

\begin{solution}{exo:g3:problems:11}
أجوبة صحيحة كثيرة --- فمثلًا: «تصنع بائعة زهور $8$ باقات
من $15$ زهرة، ثم تنزع $20$ زهرة ذابلة. فكم زهرة تبقى؟»
والحلّ: $8 \times 15 = 120$، ثم $120 - 20 = 100$.
\end{solution}
